% Keep the current LaTeX output and document hooks: tCON2e v5.0 predates
% the hook system and otherwise bypasses modern natbib/hyperref cleanup.
\makeatletter
\let\TConKernelOutputPage\@outputpage
\let\TConKernelEndDocument\enddocument
\makeatother
\documentclass{tCON2e}
\makeatletter
\let\@outputpage\TConKernelOutputPage
\let\enddocument\TConKernelEndDocument
% This is a manuscript, so omit the class's unassigned publication fields.
\let\ps@title\ps@plain
\makeatother

\usepackage[T1]{fontenc}
\usepackage{lmodern}
\usepackage{anyfontsize}
\usepackage{mathtools,bm,mathrsfs}
\usepackage{booktabs,array,enumitem}
\usepackage{needspace}
\usepackage{float}
\usepackage{xcolor}
\usepackage{microtype}
\usepackage{graphicx}
\usepackage{tikz}
\usepackage{comment}
\usetikzlibrary{arrows.meta,positioning,shapes.geometric}
\usepackage[longnamesfirst,sort]{natbib}
\bibpunct[, ]{(}{)}{;}{a}{,}{,}
\renewcommand\bibfont{\fontsize{10}{12}\selectfont}
\usepackage{hyperref}
\usepackage[nameinlink,noabbrev]{cleveref}
\newcommand{\needrev}[1]{#1}
\newenvironment{revision}{\begingroup\color{black}}{\endgroup}
\tolerance=1800
\emergencystretch=2em
\crefname{theorem}{theorem}{theorems}
\Crefname{theorem}{Theorem}{Theorems}
\crefname{proposition}{proposition}{propositions}
\Crefname{proposition}{Proposition}{Propositions}
\crefname{lemma}{lemma}{lemmas}
\Crefname{lemma}{Lemma}{Lemmas}
\crefname{corollary}{corollary}{corollaries}
\Crefname{corollary}{Corollary}{Corollaries}
\crefname{definition}{definition}{definitions}
\Crefname{definition}{Definition}{Definitions}
\crefname{remark}{remark}{remarks}
\Crefname{remark}{Remark}{Remarks}
\crefname{equation}{equation}{equations}
\Crefname{equation}{Equation}{Equations}
\crefname{figure}{Figure}{Figures}
\Crefname{figure}{Figure}{Figures}
\hypersetup{
  hidelinks,
  pdftitle={Constancy of the Wong Matrix in State Dimension Three and rank Two},
  pdfauthor={Songlin Zhou},
  pdfsubject={Finite-dimensional estimation algebras and the Wong matrix}
}


\setlength{\parindent}{1.5em}
\setlength{\parskip}{0.18em}
\setlist[itemize]{leftmargin=2.2em,itemsep=0.2em,topsep=0.3em}
\setlist[enumerate]{leftmargin=2.4em,itemsep=0.2em,topsep=0.3em}
\allowdisplaybreaks[3]
\numberwithin{equation}{section}

\newtheoremstyle{journalplain}
  {0.75em}{0.75em}{\itshape}{}{\bfseries}{.}{0.55em}{}
\newtheoremstyle{journaldef}
  {0.75em}{0.75em}{\normalfont}{}{\bfseries}{.}{0.55em}{}
\theoremstyle{journalplain}
\newtheorem{theorem}{Theorem}[section]
\newtheorem{proposition}[theorem]{Proposition}
\newtheorem{lemma}[theorem]{Lemma}
\newtheorem{corollary}[theorem]{Corollary}
\theoremstyle{journaldef}
\newtheorem{definition}[theorem]{Definition}
\newtheorem{remark}[theorem]{Remark}

\DeclareMathOperator{\ad}{ad}
\DeclareMathOperator{\ord}{ord}
\DeclareMathOperator{\Span}{span}
\DeclareMathOperator{\wt}{wt}
\DeclareMathOperator{\inop}{in}
\newcommand{\R}{\mathbb R}
\newcommand{\C}{\mathbb C}
\newcommand{\N}{\mathbb N}
\newcommand{\U}{\mathcal U}
\newcommand{\dd}{\,\mathrm d}
\newcommand{\pd}{\partial}
\newcommand{\eps}{\varepsilon}
\newcommand{\Lie}[2]{\left[#1,#2\right]}
\newcommand{\zetaSym}{\zeta}
\makeatletter\let\revision\relax\let\endrevision\relax\makeatother
\newenvironment{revision}{}{}

\articletype{RESEARCH ARTICLE}
\title{Finite Dimensional Estimation Algebra with State Dimension 3 and Rank 2: Constancy of Wong Matrix}
\author{\name{Songlin Zhou}}

\begin{document}
\maketitle
\begin{revision}
\begin{abstract}
We study the structure of finite-dimensional estimation algebras with state dimension three and rank two arising from smooth nonlinear
filtering systems, and establish the constancy of the Wong $\Omega$-matrix in this case.
Building on the linear structural results of~\citep{Shi2017Finite}, our proof combines Euler-operator techniques, and methods from polynomial dynamical systems. The central idea is to show that any nonzero
slope in the affine Wong matrix forces an infinite family of linearly
independent elements in the estimation algebra, contradicting finite
dimensionality. Consequently, the drift admits a Yau-type representation
as the sum of an affine vector field and the gradient of a smooth
function. This gives a further structural resolution of Mitter's conjecture in this setting.\footnote{The main theorem was basically proved by AI Mathematician~\citep{liu2026aim}, a mathematical proving agent,
and its proof was formally verified in Lean, which is publicly
available at \url{https://github.com/Pinarinith/AI-Mathematician/tree/main/Mitter_Conjecture/lean}. The author takes full responsibility for
the correctness of the proof and the entire content of this paper.}
\end{abstract}

\begin{keywords}
nonlinear filtering; estimation algebra; Wong matrix; non-maximal rank;
Euler operator
\end{keywords}
\begin{classcode}
17B66; 93E11; 35K15
\end{classcode}

\section{Introduction and main result}

Since the development of the Kalman-Bucy filter for linear Gaussian
systems in the early 1960s, the search for exact finite-dimensional
filters for nonlinear systems has attracted sustained interest.
The aim is to describe the conditional distribution of the state
exactly through finitely many recursively evolving statistics.
The innovations approach and the Duncan-Mortensen-Zakai (DMZ)
equation provide general formulations of nonlinear filtering,
but they do not in general yield a closed finite-dimensional
system~\citep{Fujisaki1972Stochastic,Marcus1984Algebraic}.
Numerical solution of the resulting equations typically requires
approximating the evolving conditional density or its representation
\citep{Luo2013Complete}.
The estimation-algebra approach, introduced by Brockett~\citep{Brockett1980The}, Clark~\citep{Brockett1981Nonlinear},
and Mitter~\citep{Mitter1979On},
uses the Lie algebra generated by the filtering operators to
investigate when an exact finite-dimensional realization is possible.
When this algebra is finite dimensional,
the Wei-Norman method can then express the explicit evolution
through finitely many ordinary differential equations using Lie algebraic method for time-varying differential equations
\citep{Wei1964On,Marcus1984Algebraic}.
The principal advantage of this approach is that it identifies
an algebraic mechanism for exact finite-dimensional closure,
providing a systematic framework for constructing recursive
filters. The study of estimation algebras is therefore of fundamental importance
to the theory of finite-dimensional nonlinear filtering.

Following the nonlinear filtering Lie algebra approaches of Bene\v{s} and Wong~\citep{Benes1981Exact,Wong1983New}, the classification of finite-dimensional
estimation algebras of maximal rank was completed by Yau and his
collaborators~\citep{Yau2003Complete,Yau2005Classification}.
For nonmaximal rank, Wu and Yau classified
the case of state dimension two~\citep{Wu2006Classification}. In state dimension three and linear
rank two, Shi and Yau first proved that the Wong matrix is affine
\citep{Shi2017Finite} and subsequently established Mitter's conjecture proposed in ICM 1983,
namely that every function element of the estimation algebra is affine
\citep{Shi2020Finite}\footnote{Since the main structural results of~\citep{Shi2020Finite}
provide a key foundation for our proof, we have also formalized
and verified these results in Lean, within the same verification
package as our main theorem; see
\url{https://github.com/Pinarinith/AI-Mathematician/tree/main/Mitter_Conjecture/lean}.}. More recent work has substantially extended this
picture. Jiao and Yau studied the exact subclass, in which the drift is
a gradient and the Wong matrix vanishes, combining structure results
in state dimension three and rank two with Riccati and Schr\"odinger
analysis of the associated potentials~\citep{Jiao2023Structure}.
They also extended the affine Wong-matrix theorem to arbitrary state
dimension $n$ and rank $n-1$~\citep{Jiao2024FiniteDimensional}.
Under the additional assumption of a constant Wong matrix, Yu, Jiao,
and Yau obtained a complete classification for rank $n-1$
\citep{Yu2024Complete}.

The nonmaximal-rank case is more difficult because the algebra
contains fewer independent coordinate functions, limiting the
control that commutator identities provide over the Wong matrix.
Indeed, Jiao and Yau constructed rank-$n-2$ examples in which
the Wong matrix has nonpolynomial smooth elements
\citep{Jiao2020New}.
In this paper, we combine techniques of Euler operators and methods from polynomial dynamical systems. We prove that the Wong
matrix is constant when the state dimension is three and the
linear rank is two, thereby removing the constant-Wong-matrix
assumption from the classification of \citep{Yu2024Complete}
in this case. The following is our main theorem.
\begin{theorem}[Main theorem]\label{thm:main}
Let $E$ be a finite-dimensional estimation algebra of a smooth nonlinear
filtering system on $\mathbb R^3$, as in
\eqref{eq:D-eta-L0}--\eqref{eq:estimation-algebra}, with state dimension 3 and rank 2. Then the Wong matrix $\Omega$ is constant, i.e., $\omega_{12},\omega_{13},\omega_{23}$ are constants.
\end{theorem}

The paper is organized as follows: Section~\ref{sec:2} describes some basic concepts of estimation algebra and well-known results for our proof. In Section~\ref{sec:3}, we complete the proof within 3 steps: 
\begin{enumerate}
    \item We prove that $\omega_{12}$ is constant.
    \item We prove that $\omega_{13}$ and $\omega_{23}$ do not have $x_3$-coordinate.
    \item We prove that $\omega_{13}$ and $\omega_{23}$ are constant.
\end{enumerate}

\section{Basic concepts and results.}
\label{sec:2}


\subsection{Basic concepts}
The filtering problem we consider in this paper is based on the following model.
\begin{equation}
\begin{aligned}
\mathrm dx(t)&=f(x(t))\,\mathrm dt+g(x(t))\,\mathrm dv(t),
&x(0)&=x_0,\\
\mathrm dy(t)&=h(x(t))\,\mathrm dt+\mathrm dw(t),
&y(0)&=0.
\end{aligned}
\label{eq:filtering-model}
\end{equation}
Here $x,v,y,w$ belong to $\mathbb R^n, \mathbb R^p, \mathbb R^m, \mathbb R^m$, respectively, where $x(t)$ is the state process and $y(t)$ is the observation process.
We assume that $f$ and $h=(q_1,\ldots,q_m)$ are $C^\infty$ smooth and
$g(x)g(x)^T=I_n$.
Let $\sigma$ denote the
normalized conditional density given the observation $\{y(s): 0 \le s \le t\}$, then $\sigma$ satisfies the well-known DMZ equation
\[
\mathrm d\sigma(t,x)=L_0\sigma(t,x)\mathrm dt+
\sum_{j=1}^m L_j\sigma(t,x)\mathrm dy_j, \quad \sigma(0,x) = \sigma_0,
\]
where $L_j\phi=h_j\phi$ for $1\le j\le m$ and
\[
L_0=\frac12\Delta-\sum_{i=1}^n f_i\partial_i
-\sum_{i=1}^n\partial_i f_i-\frac12\sum_{j=1}^m h_j^2.
\]
Equivalently, let
\begin{equation}
D_i=\partial_i-f_i,\quad
\eta=\sum_{i=1}^n\!\left(\partial_i f_i+f_i^2\right)+\sum_{j=1}^m h_j^2,\quad
L_0=\frac12\sum_{i=1}^n D_i^2-\frac12\eta.
\label{eq:D-eta-L0}
\end{equation}
In the following, we briefly introduce Lie algebra for differential operators and the estimation algebra.
\begin{definition}
 For two differential operators $X$ and $Y$, the Lie bracket $[,]:\mathcal F \times \mathcal F \to \mathcal F, [X,Y]:=XY-YX$ satisfying the following three axioms:
 \begin{enumerate}
     \item The Lie bracket is bilinear;
     \item $[X,X] = 0$ for all $X \in \mathcal F$;
     \item Jacobi identity: $[X,[Y,Z]] + [Y,[Z,X]] + [Z,[X,Y]] = 0$ for all $X,Y,Z \in \mathcal F$.
 \end{enumerate}
\end{definition}
\begin{definition}
The estimation algebra $E$ of \eqref{eq:filtering-model} is the smallest
real Lie algebra containing $L_0,L_1,\ldots,L_m$:
\begin{equation}
E=\operatorname{Lie}_{\mathbb R}\langle L_0,h_1,\ldots,h_m\rangle.
\label{eq:estimation-algebra}
\end{equation}
Its linear rank is $\dim_{\mathbb R}L(E)$, where
\(
L(E)=E\cap\operatorname{span}_{\mathbb R}\{x_1,x_2,...,x_n\}.
\)
A function element of $E$ is a smooth function whose multiplication
operator belongs to $E$.
\end{definition}

\begin{definition}
The Wong matrix is $\Omega=(\omega_{ij})_{1\le i,j\le n}$, with
\begin{equation}
\omega_{ij}=\partial_i f_j-\partial_j f_i=[D_j,D_i].
\label{eq:wong-def}
\end{equation}
In particular, $\omega_{ji}=-\omega_{ij}$.
\end{definition}

\begin{definition}
Let $\mathcal U_r$ be the vector space of smooth-coefficient differential
operators of order at most $r$ with finite index $I_P \subset \mathbb N^n$. $P \in \mathcal U_r$ can be represented as
\[
P=\sum_{(i_1,i_2,...,i_n)) \in I_P}a_{i_1,i_2,...,i_n}(x)\partial_1^{i_1}\partial_2^{i_2}...\partial_n^{i_n},
\]
where $|i|:=\sum_{k=1}^n i_k \le r$. The order of $P$ is denoted by $ord(P) = \max_i |i|$.
\end{definition}

Remarkably, in this paper,we write $\ad_P Q=[P,Q]$, and $ad_P^l Q = [P,ad_P^{l-1}Q]$ for iterated commutators.
\begin{comment}
the one-variable Euler operator is $t\partial_t$.
\end{comment}
\subsection{Preliminaries}
From now on, we consider the finite-dimensional algebra with state dimension $n = 3$ and linear rank 2. We choose orthonormal coordinates so that there exist $c_i$ such that $x_1 + c_i \in E$, and for any constant $c$, $x_3 + c \notin E$.
We recall the known results and methodologies needed in the proof.

\begin{theorem}[Ocone \citep{Ocone1981Finite}]\label{thm:ocone}
If $E$ is a finite-dimensional estimation algebra and a multiplication
operator $u$ belongs to $E$, then $u$ is a polynomial of degree at most two.
\end{theorem}
\begin{theorem}[See~\citet{Yau2005Classification,Wu2006Classification}]\label{thm:top-coeff}
Let $E$ be a finite-dimensional estimation algebra. If $P\in E$ has
positive differential order $r$, every coefficient of its ordinary
principal symbol $\sigma_r(P)$ is polynomial in the state variables.
Equivalently, its highest-order $D$-normal coefficients are polynomials.
\end{theorem}

\begin{theorem}[See~Theorems 3.4 and 3.10 of~\citet{Shi2017Finite}]\label{thm:shi-yau-affine}
Suppose that $E$ is finite-dimensional, the state dimension is three,
and the rank is two, then
\begin{align}
\omega_{12} &= b_1x_1+b_2x_2+b_0, \\
\omega_{13} &= k_1x_1+k_2x_2+k_3x_3+h_0, \\
\omega_{23} &= h_1x_1+h_2x_2+h_3x_3+k_0,
\end{align}
where all coefficients are real constants.
\end{theorem}

\begin{theorem}\label{thm:bianchi-affine}
For the affine Wong matrix of Theorem~\ref{thm:shi-yau-affine}, $h_1=k_2$.
\end{theorem}

\begin{proof}
Since $\omega_{ij}=\partial_if_j-\partial_jf_i$, the Bianchi identity
$\partial_3\omega_{12}-\partial_2\omega_{13}+\partial_1\omega_{23}=0$ holds
by equality of mixed partials. Substituting the affine expressions of
Theorem~\ref{thm:shi-yau-affine} gives
$\partial_3\omega_{12}=0$, $\partial_2\omega_{13}=k_2$ and
$\partial_1\omega_{23}=h_1$, so the identity reads $-k_2+h_1=0$.
\end{proof}

\begin{lemma}[Gauge conjugation]\label{prop:gauge}
Let $\Lambda\in C^\infty(\mathbb R^3;\mathbb R)$. The map
\(
\mathcal G_\Lambda(P)=e^\Lambda P e^{-\Lambda}
\)
is an automorphism of the Lie algebra of differential operators.
It fixes multiplication operators and preserves differential
order and elements of Wong matrix. In particular,
\(
E_\Lambda:=\mathcal G_\Lambda(E)
\)
is isomorphic to $E$.
\end{lemma}

\begin{proof}
Conjugation by the invertible operator $e^\Lambda$ is an algebra
automorphism, with inverse $\mathcal G_{-\Lambda}$, hence preserves the
commutator and so is a Lie algebra automorphism. If $u$ denotes
multiplication by a smooth function, then $e^\Lambda u e^{-\Lambda}=u$,
so multiplication operators are fixed. For the covariant derivative,
$e^\Lambda D_i e^{-\Lambda}=D_i+\partial_i\Lambda$, which is again a
first-order operator; conjugating the Leibniz expansion of a general
$P$ shows that the principal symbol is unchanged, so the differential
order is preserved. Finally, replacing $f$ by $f-\nabla\Lambda$ leaves
$\omega_{ij}=\partial_if_j-\partial_jf_i$ unchanged, since
$\partial_i\partial_j\Lambda-\partial_j\partial_i\Lambda=0$. Hence
$E_\Lambda\cong E$.
\end{proof}

\begin{lemma}\label{lem:polynomial-drift}
Let $f\in C^\infty(\mathbb R^3;\mathbb R^3)$, then there exists
$\Lambda\in C^\infty(\mathbb R^3;\mathbb R)$ such that
$\widetilde f=f+\nabla\Lambda$ satisfies
\begin{align}
\widetilde f_1&=0,\notag\\
\widetilde f_2
&=\frac{b_1}{2}x_1^2+b_2x_1x_2+b_0x_1,
\label{eq:f2-general}\\
\widetilde f_3
&=\frac{k_1}{2}x_1^2+k_2x_1x_2+k_3x_1x_3+k_0x_1
+\frac{h_2}{2}x_2^2+h_3x_2x_3+h_0x_2.
\label{eq:f3-general}
\end{align}
\end{lemma}

\begin{proof}
Let $\widetilde f$ be the polynomial vector field displayed above.
Direct differentiation gives
\(
\partial_1\widetilde f_2-\partial_2\widetilde f_1=\omega_{12},
\partial_1\widetilde f_3-\partial_3\widetilde f_1=\omega_{13},
\partial_2\widetilde f_3-\partial_3\widetilde f_2=\omega_{23}.
\)
Consequently, the one-form
\(
\sum_{i=1}^{3}(\widetilde f_i-f_i)\,dx_i
\)
is closed. By the Poincar\'e lemma on $\mathbb R^3$, it equals
$d\Lambda$ for some globally defined smooth real function
$\Lambda$. Hence $\widetilde f=f+\nabla\Lambda$, and
\(
e^\Lambda D_i e^{-\Lambda}
=
\partial_i-f_i-\partial_i\Lambda
=
\widetilde D_i.
\)
\end{proof}


\begin{lemma}[\citep{Wu2006Classification}]\label{lem:operator-products-symbols}
Let $g,h\in C^\infty(\mathbb R^n)$ and
$\alpha,\beta\in\mathbb N_0^n$, put $r=|\alpha|$ and $s=|\beta|$. If $r+s\geq2$, then
\[
[gD^\alpha,hD^\beta]
\equiv
\sum_{\substack{1\leq\ell\leq n\\
                 \alpha_\ell+\beta_\ell>0}}
\left(
\alpha_\ell g\,\partial_\ell h
-\beta_\ell h\,\partial_\ell g
\right)
D^{\alpha+\beta-e_\ell}
\pmod{\mathcal U_{r+s-2}},
\]
where $e_\ell$ is the $\ell$th standard multi-index.
\end{lemma}

\begin{lemma}[\citep{Yau1994FiniteDimensional,Chiou1994FiniteDimensional}]\label{lem:commutator-identities}
Let $E$ be an estimation algebra, let $X,Y,Z\in E$, and let
$g,h\in C^\infty(\mathbb R^n)$. With the notation above, the
following identities hold for $1\leq i,j,k\leq n$:
\begin{enumerate}
\renewcommand{\labelenumi}{(\arabic{enumi})}

\item
\(
[XY,Z]=X[Y,Z]+[X,Z]Y.
\)


\item
\(
\displaystyle
[L_0,g]
=
\sum_{a=1}^{n}\frac{\partial g}{\partial x_a}D_a
+\frac12\Delta g.
\)

\item
\(
\displaystyle
[L_0,D_i]
=
\sum_{a=1}^{n}\omega_{ia}D_a
+\frac12\sum_{a=1}^{n}
\frac{\partial\omega_{ia}}{\partial x_a}
+\frac12\frac{\partial\eta}{\partial x_i}.
\)


\item
\(
\displaystyle
[gD_i,hD_j]
=
gh\omega_{ji}
+g\frac{\partial h}{\partial x_i}D_j
-h\frac{\partial g}{\partial x_j}D_i,
\)
where $\omega_{ji}=[D_i,D_j]$.

\item
\(
\displaystyle
[gD_i^2,h]
=
2g\frac{\partial h}{\partial x_i}D_i
+g\frac{\partial^2h}{\partial x_i^2}.
\)

\item
\(
\begin{aligned}[t]
[D_i^2,hD_j]
&=
2\frac{\partial h}{\partial x_i}D_iD_j
+2h\omega_{ji}D_i
+\frac{\partial^2h}{\partial x_i^2}D_j
+h\frac{\partial\omega_{ji}}{\partial x_i}.
\end{aligned}
\)

\item
\(
\begin{aligned}[t]
[D_i^2,D_j^2]
&=
4\omega_{ji}D_jD_i
+2\frac{\partial\omega_{ji}}{\partial x_j}D_i
+2\frac{\partial\omega_{ji}}{\partial x_i}D_j
+\frac{\partial^2\omega_{ji}}{\partial x_i\partial x_j}
+2\omega_{ji}^2.
\end{aligned}
\)

\item
\(
\begin{aligned}[t]
[D_k^2,hD_iD_j]
&=
2\frac{\partial h}{\partial x_k}D_kD_iD_j
+2h\omega_{jk}D_iD_k
+2h\omega_{ik}D_kD_j
+\frac{\partial^2h}{\partial x_k^2}D_iD_j\\
&+2h\frac{\partial\omega_{jk}}{\partial x_i}D_k
+h\frac{\partial\omega_{jk}}{\partial x_k}D_i
+h\frac{\partial\omega_{ik}}{\partial x_k}D_j
+h\frac{\partial^2\omega_{jk}}{\partial x_i\partial x_k}.
\end{aligned}
\)

\item
\(
\begin{aligned}[t]
[gD_iD_j,hD_k]
&=
g\frac{\partial h}{\partial x_j}D_iD_k
+g\frac{\partial h}{\partial x_i}D_jD_k
-h\frac{\partial g}{\partial x_k}D_iD_j\\
&\quad
+gh\omega_{kj}D_i
+gh\omega_{ki}D_j
+g\frac{\partial^2h}{\partial x_i\partial x_j}D_k
+gh\frac{\partial\omega_{kj}}{\partial x_i}.
\end{aligned}
\)

\end{enumerate}
\end{lemma}

\section{Proof of the main theorem}\label{sec:3}
\subsection{$\omega_{12}$ is constant}\label{sec:visible}

First, we consider
\begin{align}
H_0
&=[L_0,D_1]-\frac12(b_2+k_3)
 =\omega_{12}D_2+\omega_{13}D_3+\frac12\eta_1 \in E,
\label{eq:H0bar}\\
H_1
&=[L_0,D_2]-\frac12(-b_1+h_3)
 =-\omega_{12}D_1+\omega_{23}D_3+\frac12\eta_2 \in E.
\label{eq:H1bar}
\end{align}
For $B\in E\cap\mathcal U_2$, write
\[
B=a_{11}D_1^2+a_{12}D_1D_2+a_{22}D_2^2+\cdots
\pmod{\mathcal U_1},
\]
where the omitted second-order terms contain $D_3$.
Direct computation gives
\(
[[B,x_1],x_1]=2a_{11},
[[B,x_1],x_2]=a_{12},
[[B,x_2],x_2]=2a_{22}.
\)
These functions belong to $E$. By the preceding result on function
elements,
\(
a_{11},a_{12},a_{22}
\in\operatorname{span}_{\mathbb R}\{1,x_1,x_2\}.
\)

\begin{lemma}\label{lem:linear-diagonal-coefficient}
Let $C\in E\cap\mathcal U_2$. If, for $i=1$ or $i=2$,
\[
C=(Ax_i+c)D_i^2+\cdots\pmod{\mathcal U_1},
\]
where $A,c\in\mathbb R$, then $A=0$.
\end{lemma}

\begin{proof}
It suffices to consider $i=1$. Define
\(
P_0=[L_0,C], P_{r+1}=[P_r,C]\quad(r\geq0).
\)
Each $P_r$ belongs to $E$. The commutator formulas give
\[
P_0=AD_1^3+\cdots\pmod{\mathcal U_2}
\]
and, inductively,
\[
P_r=\frac{(r+2)!}{2}A^{r+1}D_1^{r+3}
     +\cdots\pmod{\mathcal U_{r+2}}.
\]
Here the terms containing $D_2$ or $D_3$ with order $r+3$ are omitted.
Indeed, since
\(
[D_1^{r+3},C]
=(r+3)A D_1^{r+4}+\cdots\pmod{\mathcal U_{r+3}},
\)
commuting two omitted terms leaves at least one $D_2$ or $D_3$
factor at order $r+4$. Commuting an omitted term with a displayed
term can remove its last such factor only by differentiating
$Ax_1+c$, or the constant displayed coefficient of $P_r$, in
$x_2$ or $x_3$; these derivatives vanish. The terms arising from
$[D_i,D_j]=\omega_{ji}$ have lower differential order and are
included in $\mathcal U_{r+3}$.

If $A\ne0$, $P_r$ has order $r+3$ and lies in $E$, contradicting finite-dimensionality
of $E$. The case $i=2$ is identical after interchanging $x_1$ and $x_2$.
\end{proof}

\begin{lemma}\label{lem:constant-quadratic-coefficients}
Under the assumptions of the main theorem, suppose that
$b_1=0$ and $b_2\ne0$, so that
\[
\omega_{12}=b_2x_2+b_0.
\]
Then, for every $B\in E\cap\mathcal U_2$, the coefficients of $D_1D_2$
and $D_2^2$ are constant.
\end{lemma}

\begin{proof}
Direct calculation gives
\begin{align*}
[H_0,B]
&=\omega_{12}(\partial_2a_{11})D_1^2
 +(\omega_{12}\partial_2a_{12}-b_2a_{12})D_1D_2
 +(\omega_{12}\partial_2a_{22}-2b_2a_{22})D_2^2
 +\cdots\pmod{\mathcal U_1},\\
[H_1,B]
&=(-\omega_{12}\partial_1a_{11}+b_2a_{12})D_1^2
 +(-\omega_{12}\partial_1a_{12}+2b_2a_{22})D_1D_2
 -\omega_{12}(\partial_1a_{22})D_2^2
 +\cdots\pmod{\mathcal U_1}.
\end{align*}
The omitted second-order terms again contain $D_3$.
Terms involving $D_3$ in $B$ do not contribute to the three
displayed coefficients, because
$\partial_3\omega_{12}=\partial_3a_{11}
=\partial_3a_{12}=\partial_3a_{22}=0$.
Both commutators belong to $E\cap\mathcal U_2$.

Since $a_{11}$ is affine,
\[
\begin{aligned}
{}[H_0,B]-b_2B
={}&\bigl(-b_2(\partial_1a_{11})x_1
       +b_0\partial_2a_{11}-b_2a_{11}(0)\bigr)D_1^2+\cdots\pmod{\mathcal U_1},
\end{aligned}
\]
where the omitted terms are the other second-order terms.
Lemma~\ref{lem:linear-diagonal-coefficient} gives
$\partial_1a_{11}=0$. This holds for every
$B\in E\cap\mathcal U_2$.

Similarly,
\(
[H_1,B]
=-(\partial_1a_{22})(b_2x_2+b_0)D_2^2
 +\cdots\pmod{\mathcal U_1},
\)
where the omitted terms are the other second-order terms.
The same lemma gives $\partial_1a_{22}=0$.
Applying it once more to $B$ gives $\partial_2a_{22}=0$.
Thus $a_{22}$ is constant.

The coefficient of $D_1^2$ in $[H_1,B]$ is now $b_2a_{12}$.
By the conclusion already obtained for every element of
$E\cap\mathcal U_2$, this coefficient is independent of $x_1$.
Hence $\partial_1a_{12}=0$.
It remains to prove $\partial_2a_{12}=0$.

Define
\(
P_0=[L_0,B],P_{r+1}=[P_r,B]
\)
Since $a_{11},a_{12}$ are affine functions of $x_2$ and $a_{22}$
is constant, direct calculation gives
\[
P_0=(\partial_2a_{11})D_1^2D_2
    +(\partial_2a_{12})D_1D_2^2
    +\cdots\pmod{\mathcal U_2}.
\]
Iteratively repeated application of the same commutator formulas yields,
for $r\geq1$,
\[
\begin{aligned}
P_r
={}&2^r(\partial_2a_{12})^{r+1}D_1^{r+1}D_2^2\\
&+(2^{r+1}-1)(\partial_2a_{11})(\partial_2a_{12})^r
  D_1^{r+2}D_2\\
&+(2^r-1)(\partial_2a_{11})^2(\partial_2a_{12})^{r-1}
  D_1^{r+3}
 +\cdots\pmod{\mathcal U_{r+2}}.
\end{aligned}
\]
To check the induction, commute each of
the three displayed monomials with $B$. Only differentiation of
$a_{11}$ and $a_{12}$ with respect to $x_2$ contributes to terms
of the next order that do not contain $D_3$. The omitted terms
cannot contribute such terms, because all displayed coefficients
are independent of $x_3$.

If $\partial_2a_{12}\ne0$, the coefficient of
$D_1^{r+1}D_2^2$ is nonzero for every $r$, so the orders of the
elements $P_r\in E$ are unbounded. This contradicts
finite-dimensionality.
$a_{12}$ is constant as well.
\end{proof}

\begin{theorem}\label{thm:w12-constant}
Under the assumptions of the main theorem,
$b_1=b_2=0$, so $\omega_{12}=b_0$.
\end{theorem}

\begin{proof}
Suppose that $(b_1,b_2)\ne(0,0)$. Make the orthogonal change
of coordinates
\[
\begin{pmatrix}x_1\\x_2\end{pmatrix}
=
\frac{1}{\sqrt{b_1^2+b_2^2}}
\begin{pmatrix}b_2&b_1\\-b_1&b_2\end{pmatrix}
\begin{pmatrix}y_1\\y_2\end{pmatrix},
\qquad x_3=y_3.
\]
This transformation preserves the second-order
part of $L_0$ and transforms $\omega_{12}$ into
$\sqrt{b_1^2+b_2^2}\,y_2+b_0$.
The transformed algebra is isomorphic to $E$, and $\omega_{13},\omega_{23}$ stay affine.
After renaming the coordinates and coefficients, we may
therefore assume
\(
b_1=0,b_2>0,
\omega_{12}=b_2x_2+b_0.
\)
Direct calculation gives
\[
[D_1,H_0]=k_1D_3\pmod{\mathcal U_0},
\qquad
[D_1,H_1]=k_2D_3\pmod{\mathcal U_0}.
\]
Since $[D_3,x_1]=[D_3,x_2]=0$ and
$\partial_3\omega_{12}=0$, it follows that
\begin{equation}
\begin{aligned}
{}[[D_1,H_0],x_i]&=[[D_1,H_1],x_i]=0
\qquad(i=1,2),\\
[\omega_{12},[D_1,H_1]]&=0.
\end{aligned}
\label{eq:jacobi-vertical}
\end{equation}

The constant terms subtracted in
\eqref{eq:H0bar}--\eqref{eq:H1bar} do not affect the following
commutators. By the Jacobi identity,
\[
[D_2,H_0]-[D_1,H_1]
=[L_0,\omega_{12}]=b_2D_2.
\]
Using $[D_2,D_1]=\omega_{12}$ and
$[H_0,\omega_{12}]=b_2\omega_{12}$, we obtain
\begin{equation}
\begin{aligned}
{}[D_2,[D_1,H_0]]-[D_1,[D_1,H_1]]
&=[D_1,b_2D_2]+[\omega_{12},H_0]=-2b_2\omega_{12}.
\end{aligned}
\label{eq:jacobi-first-obstruction}
\end{equation}

Next, direct calculation gives
\[
[L_0,H_0]
=b_2D_2^2+k_1D_1D_3+k_2D_2D_3+k_3D_3^2
\pmod{\mathcal U_1}.
\]
All displayed coefficients are constant. Since
\(
[L_0,D_iD_j]
=[L_0,D_i]D_j+D_i[L_0,D_j]
\in\mathcal U_2,
\)
and the commutator of $L_0$ with a first-order operator
has order at most two, we have
\(
[L_0,[L_0,H_0]]\in E\cap\mathcal U_2.
\)

Using $[H_0,x_1]=0$, $[H_0,x_2]=\omega_{12}$, and the
Jacobi identity, we obtain
\begin{align*}
[[L_0,H_0],x_1]&=[D_1,H_0],\\
[[L_0,H_0],x_2]&=[D_1,H_1]+2b_2D_2.\\
\bigl[[[L_0,[L_0,H_0]],x_1],x_2\bigr]
&=2[D_2,[D_1,H_0]],\\
\bigl[[[L_0,[L_0,H_0]],x_2],x_2\bigr]
&=2[D_2,[D_1,H_1]].
\label{eq:jacobi-head-extraction}
\end{align*}
The left-hand sides of the last two equations are, respectively, the coefficient
of $D_1D_2$ and twice the coefficient of $D_2^2$ in
$[L_0,[L_0,H_0]] \in E\cap\mathcal U_2$.
By Lemma~\ref{lem:constant-quadratic-coefficients},
these coefficients are constant. Consequently,
\[
[D_2,[D_1,H_0]]
\quad\text{and}\quad
[D_2,[D_1,H_1]]
\]
are constant multiplication operators.

Finally, take the commutator of both sides of
\eqref{eq:jacobi-first-obstruction} with $D_2$:
\[
\begin{aligned}
{}[D_2,[D_1,[D_1,H_1]]]
&=[D_1,[D_2,[D_1,H_1]]]
  +[\omega_{12},[D_1,H_1]]=[D_1,[D_2,[D_1,H_1]]],
\end{aligned}
\]
where the last equality follows from
\eqref{eq:jacobi-vertical}.
The constant multiplication operators above commute
with $D_1$ and $D_2$. Therefore,
\[
\begin{aligned}
0=[D_2,[D_2,[D_1,H_0]]]
  -[D_1,[D_2,[D_1,H_1]]]=[D_2,-2b_2\omega_{12}]
=-2b_2^2\,1.
\end{aligned}
\]
This contradicts $b_2>0$ and proves the theorem.
\end{proof}


\subsection{$h_3$ = $k_3$ = 0.}

By Theorem~\ref{thm:w12-constant}, $b_1=b_2=0$.
Suppose that $\mu=k_3^2+h_3^2>0$.
By Lemmas~\ref{lem:polynomial-drift} and~\ref{prop:gauge},
we may use the drift in
\eqref{eq:f2-general}--\eqref{eq:f3-general}, retaining
the notation $E,D_i,L_0,H_i$. In particular,
$f_1=0$ and $f_2=b_0x_1$.

Choose the affine coordinate
\[
t=\frac{k_3\omega_{13}+h_3\omega_{23}}{\mu}
=x_3+\alpha x_1+\beta x_2+\delta,
\]
where
\(
\alpha=\frac{k_3k_1+h_3k_2}{\mu},
\beta=\frac{k_3k_2+h_3h_2}{\mu},
\delta=\frac{k_3k_0+h_3h_0}{\mu}.
\)
The coordinate derivatives in $(x_1,x_2,t)$ are
\(
\nabla_1=\partial_1-\alpha\partial_3,
\nabla_2=\partial_2-\beta\partial_3,
\partial_t=\partial_3,
\)
where $\partial_i$ denotes differentiation in the original
coordinates. Then, the Euler operator
\[
J:=\frac{k_3(H_0-b_0D_2)+h_3(H_1+b_0D_1)}{\mu}
=t\partial_3+r(x_1,x_2,t)\in E
\]
for a smooth function $r$.
Take
\[
\Lambda(x_1,x_2,t)
=\int_0^t
\frac{r(x_1,x_2,s)-r(x_1,x_2,0)}{s}\,\mathrm ds,
\]
where the integrand extends smoothly across $s=0$.
Writing $\widehat P=e^\Lambda Pe^{-\Lambda}$ and
$\widehat E=e^\Lambda Ee^{-\Lambda}$, we obtain
\(
\widehat J=t\partial_3+r(x_1,x_2,0),
\)
and
\(
\widehat D_1
=\nabla_1+\alpha\partial_3-\partial_1\Lambda,
\widehat D_2
=\nabla_2+\beta\partial_3-b_0x_1-\partial_2\Lambda.
\)

To extract the required terms involving $\partial_3$ as elements of
the estimation algebra, we establish the following lemma. It shows
that projection onto each weight component preserves membership
in the algebra for the operators in $\widehat E$.

\begin{lemma}\label{lem:smooth-euler-module}
Let $\widehat E,\widehat D_1,\widehat D_2,\widehat L_0$
be defined as above, and suppose that
\[
\widehat J=t\partial_t+r_0(x_1,x_2)\in\widehat E.
\]
For each $P\in\{\widehat D_1,\widehat D_2,\widehat L_0\}$,
write
\[
P=
\sum_{\substack{i,j,q\ge0\\i+j+q\le2}}
a_{ijq}(x_1,x_2,t)
\partial_{x_1}^i\partial_{x_2}^j\partial_t^q.
\]
(1) Then every coefficient is polynomial in $t$:
\[
a_{ijq}(x_1,x_2,t)
=\sum_{p\ge0}a_{ijqp}(x_1,x_2)t^p,\quad
a_{ijqp}\in C^\infty(\mathbb R^2).
\]
(2) Moreover, define the \textbf{weight} of a nonzero term by
\(
\operatorname{wt}\!\left(
a_{ijqp}(x_1,x_2)t^p
\partial_{x_1}^i\partial_{x_2}^j\partial_t^q
\right)=p-q.
\)
Let $P_w$ denote the sum of all terms of weight $w$ in $P$, then for every $w\in\mathbb Z$, there exists
$\pi_w\in\mathbb R[s]$ such that
\[
P_w=\pi_w(\operatorname{ad}_{\widehat J})P
\in\widehat E.
\]
\end{lemma}

\begin{proof}
Direct calculation gives
\[
[t\partial_t,a\partial_{x_1}^i\partial_{x_2}^j\partial_t^q]
=(t\partial_ta-qa)\partial_{x_1}^i\partial_{x_2}^j\partial_t^q.
\]
In a commutator with $r_0$, at least one derivative in
$x_1,x_2$ falls on $r_0$. Thus commutation with $r_0$
lowers the differential order in these variables.
Since $r_0$ is independent of $t$, so the Jacobi
identity gives
\(
[\operatorname{ad}_{t\partial_t},\operatorname{ad}_{r_0}]
=\operatorname{ad}_{[t\partial_t,r_0]}=0.
\)

Let $\chi$ be the characteristic polynomial of
$\operatorname{ad}_{\widehat J}$ on $\widehat E$.
By the Cayley-Hamilton theorem,
$\chi(\operatorname{ad}_{\widehat J})P=0$.
Since $r_0$ is a function and $P\in\mathcal U_2\cap\widehat E$,
each commutator with $r_0$ lowers the differential order,
so $(\operatorname{ad}_{r_0})^3P=0$.
Thus
\[
\begin{aligned}
0
&=\bigl[\chi(\operatorname{ad}_{\widehat J})\bigr]^3P\\
&=\bigl[\chi(\operatorname{ad}_{t\partial_t})\bigr]^3P
+3\chi(\operatorname{ad}_{t\partial_t})
\bigl[\chi(\operatorname{ad}_{\widehat J})
-\chi(\operatorname{ad}_{t\partial_t})\bigr]
\chi(\operatorname{ad}_{\widehat J})P\\
&\quad+
\left(
\chi'(\operatorname{ad}_{t\partial_t})
+\frac12\chi''(\operatorname{ad}_{t\partial_t})
\operatorname{ad}_{r_0}
\right)^3
(\operatorname{ad}_{r_0})^3P\\
&=\bigl[\chi(\operatorname{ad}_{t\partial_t})\bigr]^3P,
\end{aligned}
\]
where the expansion uses
$[\operatorname{ad}_{t\partial_t},\operatorname{ad}_{r_0}]=0$,
and the last equality follows from
$\chi(\operatorname{ad}_{\widehat J})P=0$
and $(\operatorname{ad}_{r_0})^3P=0$. Comparing their coefficients gives
\[
[\chi(t\partial_t-q)]^3a_{ijq}=0.
\]
Applying Lemma~\ref{lem:global-smooth-euler} for each fixed
$(x_1,x_2)$ yields
\[
a_{ijq}(x_1,x_2,t)
=
\sum_{\substack{p\ge0\\\chi(p-q)=0}}
\frac{\partial_t^p a_{ijq}(x_1,x_2,0)}{p!}\,t^p.
\]
There are only finitely many such $p$, independently of
$(x_1,x_2)$, and the displayed coefficients are smooth.
This proves (1).

For (2), write $P=\sum_w P_w$, a finite sum.
By the definition of weight,
\[
[t\partial_t,P_w]=wP_w.
\]
Commutation with $r_0$ changes neither the power of $t$
nor the number of $\partial_t$ factors, so it preserves
weight. It also lowers the differential order in
$x_1,x_2$. Consequently,
\[
(\operatorname{ad}_{\widehat J}-w)^3P_w
=(\operatorname{ad}_{r_0})^3P_w=0.
\]

Fix an occurring weight $w$. Since the polynomials
$(s-v)^3$ for distinct occurring weights $v$ are pairwise
coprime, the Chinese remainder theorem gives
$\pi_w\in\mathbb R[s]$ satisfying
\[
\begin{aligned}
\pi_w(s)&\equiv1\pmod{(s-w)^3},\\
\pi_w(s)&\equiv0\pmod{(s-v)^3}
\qquad\text{for every other finitely occurring weight }v.
\end{aligned}
\]
The preceding annihilation identities therefore imply
\[
\pi_w(\operatorname{ad}_{\widehat J})P
=
\sum_v\pi_w(\operatorname{ad}_{\widehat J})P_v
=P_w.
\]
Since $\widehat J\in\widehat E$, its adjoint action
preserves $\widehat E$, and hence $P_w\in\widehat E$.
For a weight that does not occur, take $\pi_w=0$.
\end{proof}

\begin{theorem}\label{thm:sector-I-smooth}
Under the assumptions of the main theorem, $k_3=h_3=0$.
\end{theorem}

\begin{proof}
By Lemma~\ref{lem:smooth-euler-module},
$\partial_1\Lambda$ and $\partial_2\Lambda$ are polynomial
in $t$, and their positive-weight components are
multiplication operators in $\widehat E$.
Since the coordinate change leaves $x_1,x_2$ fixed and
gauge conjugation fixes multiplication operators, we also have
$
\widehat E\cap C^\infty(\mathbb R^3)
=\operatorname{span}_{\mathbb R}\{1,x_1,x_2\}.$
Hence $\partial_1\Lambda$ and $\partial_2\Lambda$ are
independent of $t$, while
$
\partial_3\Lambda=\rho(x_3)
=\rho(t-\alpha x_1-\beta x_2-\delta)
$
for some $\rho\in C^\infty(\mathbb R)$.

Using \eqref{eq:f3-general}, direct expansion gives
\begin{align}
\widehat L_0
&=
\underbrace{\frac12(\nabla_1^2+\nabla_2^2)}_{\text{weight }0}
+
\underbrace{\alpha\nabla_1\partial_3
+\beta\nabla_2\partial_3}_{\text{weight }-1}
+
\underbrace{\frac{1+\alpha^2+\beta^2}{2}\partial_3^2}_{\text{weight }-2}
\notag\\
&\quad
\underbrace{+a_1\nabla_1+a_2\nabla_2
-\bigl((k_3x_1+h_3x_2)t
+\rho(t-\alpha x_1-\beta x_2-\delta)+b\bigr)\partial_3
+u(x_1,x_2,t)}_{\text{weight }\ge 0,\text{ depending on the order of } \rho}.
\label{eq:Lhat-short}
\end{align}
where $a_1,a_2,b$ depend only on $x_1,x_2$.
Lemma~\ref{lem:smooth-euler-module} gives
$u\in C^\infty(\mathbb R^2)[t]$, $\rho\in\mathbb R[t]$ is a polynomial, and
the $-2$ weight term is
\[
(\widehat L_0)_{-2}
=\frac{1+\alpha^2+\beta^2}{2}\partial_3^2\in\widehat E,
\qquad
\partial_3^2\in\widehat E.
\]

\emph{Case 1: $\deg\rho=d\ge2$.}
Write $\rho(s)=a_ds^d+\cdots$, with $a_d\ne0$.
Then
\[
(\widehat L_0)_{d-1}
=-a_dt^d\partial_3\pmod{\mathcal U_0}
\]
belongs to $\widehat E$.
Start with $P_2=\partial_3^2$. If
$P_m=\partial_3^m\pmod{\mathcal U_{m-1}}$ belongs to
$\widehat E$, the commutator formulas give
\begin{align*}
[(\widehat L_0)_{d-1},P_m]
&=m a_d d\,t^{d-1}\partial_3^m
\pmod{\mathcal U_{m-1}},\\
[P_m,[(\widehat L_0)_{d-1},P_m]]
&=m^2 a_d d(d-1)t^{d-2}\partial_3^{2m-1}
\pmod{\mathcal U_{2m-2}}.
\end{align*}
Commuting $d-2$ times with $\partial_3^2$ gives
\[
\begin{aligned}
&\operatorname{ad}_{\partial_3^2}^{\,d-2}
\bigl([P_m,[(\widehat L_0)_{d-1},P_m]]\bigr)
=2^{d-2}m^2 a_d d!\,\partial_3^{2m+d-3}
\pmod{\mathcal U_{2m+d-4}}.
\end{aligned}
\]
Dividing by the displayed nonzero constant produces
$P_{2m+d-3}$ with the same required form.
Since $2m+d-3>m$, iteration gives elements of unbounded
differential order, contradicting finite-dimensionality.

\emph{Case 2: $\rho(s)=cs+c_0$.}
This includes $\rho=0$. The weight-zero component is
\[
\begin{aligned}
(\widehat L_0)_0
&=\frac12(\nabla_1^2+\nabla_2^2)
+a_1\nabla_1+a_2\nabla_2
-(k_3x_1+h_3x_2+c)t\partial_3+u_0(x_1,x_2)
\in\widehat E.
\end{aligned}
\]
If $\partial_3^m\in\widehat E$ for $m\ge2$, then
\[
\frac{[(\widehat L_0)_0,\partial_3^m]-mc\partial_3^m}{m}
=(k_3x_1+h_3x_2)\partial_3^m\in\widehat E.
\]
A further calculation gives
\[
\begin{aligned}
&\bigl[(k_3x_1+h_3x_2)\partial_3^m,
[(\widehat L_0)_0,(k_3x_1+h_3x_2)\partial_3^m]\bigr]
=-(k_3^2+h_3^2)\partial_3^{2m}
=-\mu\partial_3^{2m} \in \widehat E.
\end{aligned}
\]
which is a contradiction against finite-dimensionality.
Therefore $k_3=h_3=0$.
\end{proof}

\section{$k_1=k_2=h_1 = h_2=0$}\label{sec:mixed}


\begin{lemma}\label{lem:constant-wronskian}
Let $p,q\in\mathbb R[\tau]$. If $pq'-p'q$ is a nonzero
constant, then
\[
\operatorname{span}_{\mathbb R}\{p,q\}
=\operatorname{span}_{\mathbb R}\{1,\tau\}.
\]
\end{lemma}

\begin{proof}
The nonzero Wronskian implies that $p,q$ are linearly
independent. If their degrees coincide, subtract a constant
multiple of one from the other to cancel its leading term.
After interchanging them if necessary, their degrees satisfy
$r<s$. Their Wronskian has degree $r+s-1$, with nonzero
leading coefficient
$(s-r)\operatorname{lc}(p)\operatorname{lc}(q)$, where $\operatorname{lc}(p)$ means the leading coefficient of polynomial $p$.
Thus $r+s=1$, so $r=0$ and $s=1$, proving the assertion.
\end{proof}

\begin{theorem}\label{thm:remaining-visible}
Under the assumptions of the main theorem, $k_1=k_2=h_1=h_2=0$.
\end{theorem}

\noindent\emph{Sketch of proof.}
We argue by contradiction. If the remaining coefficient matrix \(
S_n=(\operatorname{ad}_{\widetilde L_0})^n(\partial_1)
\in\widetilde E.
\)
is nonzero, an orthogonal coordinate change and a translation
reduce the problem to $k_1\ne0$, $k_2=k_0=0$.
We can restrict the $x_3$-dependence of $\eta$ and produces a
first-order element of the form
$\partial_3+g(x_1,x_2)+\varepsilon(x_3)$.

A gauge conjugation removes $\varepsilon$ (which is a polynomial by Theorem~\ref{thm:top-coeff}) from this first-order
element while preserving finite-dimensionality.
Further commutators, together with the function-element theorem,
show that the remaining scalar coefficient is also polynomial.
Consequently, every coefficient of $\widetilde L_0$ is polynomial
in $x_3$.

We then consider a sequence of Lie algebra elements
\[
S_n=(\operatorname{ad}_{\widetilde L_0})^n(\partial_1)
\in\widetilde E.
\]
Their coefficients remain polynomial in $x_3$, so each $S_n$
has finite "weighted-degrees". Since their span is finite-dimensional,
these weights must have a common upper bound.

We contradict this bound in two cases.
If $\deg\varepsilon\ge3$, direct commutator calculations give
terms whose nonzero highest weighted-degrees increases without bound.
If $\deg\varepsilon\le2$, these terms satisfy a polynomial recurrence
described by an associated dynamical system.
If the recurrence vanished at some stage, suitable trajectories
would yield two linearly independent polynomial solutions of
\[
u''=k_1\zeta(\tau)u,\qquad \zeta(0)=1.
\]
Their Wronskian is a nonzero constant, so
Lemma~\ref{lem:constant-wronskian} would force their span to
contain a nonzero constant solution, contradicting $k_1\ne0$.
Thus the highest-weight terms never vanish, and the weighted-degrees
again grow without bound, deriving contradictions.

\begin{proof}
By Theorems~\ref{thm:w12-constant} and~\ref{thm:sector-I-smooth},
$\omega_{12}=b_0$ and $k_3=h_3=0$. By Theorem~\ref{thm:bianchi-affine} we have $h_1 = k_2$, so it suffices to exclude the case
\[
\begin{pmatrix}k_1&k_2\\k_2&h_2\end{pmatrix}\ne0.
\]
We argue by contradiction. An orthogonal change of $(x_1,x_2)$
diagonalizes the matrix above. Choose a nonzero eigenvalue as the first
diagonal entry and translate $x_1$ to remove the constant term in
$\omega_{13}$. After renaming the coordinates and coefficients, we have
\[
k_1\ne0,\qquad k_2=k_0=0,\qquad
\omega_{13}=k_1x_1,\qquad \omega_{23}=h_2x_2+h_0.
\]
By Lemmas~\ref{lem:polynomial-drift} and~\ref{prop:gauge},
we may also assume
\(
f_1=0,f_2=b_0x_1,
f_3=\frac{k_1}{2}x_1^2+\frac{h_2}{2}x_2^2+h_0x_2,
\)
retaining the notation $E,D_i,L_0,H_i,\eta$.

\medskip
\noindent\emph{Step 1: Determine the $x_3$-dependence of $\eta$.}

Direct calculation gives
\begin{align}
[D_1,H_0]
&=k_1D_3-b_0^2-k_1^2x_1^2+\tfrac12\eta_{11},\label{eq:1}\\
[D_2,H_0]
&=-k_1x_1(h_2x_2+h_0)+\tfrac12\eta_{12},\label{eq:2}\\
[D_2,H_1]
&=h_2D_3-b_0^2-(h_2x_2+h_0)^2+\tfrac12\eta_{22},\label{eq:3}
\end{align}
and hence
\begin{align}
[D_1,[D_1,H_0]]
&=\tfrac12\partial_1^3(\eta-f_3^2),\label{eq:4}\\
[D_2,[D_1,H_0]]
&=\tfrac12\partial_1^2\partial_2(\eta-f_3^2),\label{eq:5}\\
[D_2,[D_2,H_0]]
&=\tfrac12\partial_1\partial_2^2(\eta-f_3^2),\label{eq:6}\\
[D_2,[D_2,H_1]]
&=\tfrac12\partial_2^3(\eta-f_3^2).\label{eq:7}
\end{align}
These are function elements of $E$, and therefore belong to
$\operatorname{span}_{\mathbb R}\{1,x_1,x_2\}$.

Set $\eta_0(x_1,x_2)=\eta(x_1,x_2,0)$.
Since $f_3$ is independent of $x_3$, the preceding identities~\ref{eq:4}-\ref{eq:7}
show that all third-order derivatives of $\eta-\eta_0$
with respect to $x_1,x_2$ vanish. Taylor's formula therefore gives
\[
\begin{aligned}
\eta={}&\eta_0(x_1,x_2)
+a_{11}(x_3)x_1^2+a_{12}(x_3)x_1x_2+a_{22}(x_3)x_2^2\\
&+\ell_1(x_3)x_1+\ell_2(x_3)x_2+e(x_3),
\end{aligned}
\]
where all six coefficient functions are smooth and vanish
at $x_3=0$.

Moreover,
\(
[D_2,H_0]\) and \(
h_2[D_1,H_0]-k_1[D_2,H_1]
\)
are function elements of $E$ and hence are independent of $x_3$.
The equations~\ref{eq:1}-\ref{eq:3} consequently imply that
$\eta_{12}$ and $h_2\eta_{11}-k_1\eta_{22}$ are independent of $x_3$.
Substituting the expansion above gives
\[
a_{12}'(x_3)=0,
\qquad
h_2a_{11}'(x_3)-k_1a_{22}'(x_3)=0.
\]
Since $a_{11}(0)=a_{12}(0)=a_{22}(0)=0$, it follows that
\(
a_{12}=0,
h_2a_{11}=k_1a_{22}.
\)
As $k_1\ne0$, define
\(
\varepsilon(x_3)=\frac{a_{11}(x_3)}{k_1}.
\)
Then $a_{11}=k_1\varepsilon$ and $a_{22}=h_2\varepsilon$.
Using
\(
f_3=\frac{k_1}{2}x_1^2+\frac{h_2}{2}x_2^2+h_0x_2,
\)
we obtain
\(
a_{11}x_1^2+a_{22}x_2^2
=2\varepsilon f_3-2h_0\varepsilon x_2.
\)
Replacing $\ell_2$ by $\ell_2-2h_0\varepsilon$ and retaining
the same notation, we thus have
\[
\eta=\eta_0(x_1,x_2)+2\varepsilon(x_3)f_3
+\ell_1(x_3)x_1+\ell_2(x_3)x_2+e(x_3).
\]
Here $\varepsilon,\ell_1,\ell_2,e$ are smooth and still vanish
at $x_3=0$.

The identity for $[D_1,H_0]$ now gives
\[
T=\frac1{k_1}[D_1,H_0]
=\partial_3+g(x_1,x_2)+\varepsilon(x_3)\in E
\]
for some $g\in C^\infty(\mathbb R^2)$.
Direct calculation yields
\[
[T,H_0-b_0D_2]=\tfrac12\ell_1'(x_3),
\qquad
[T,H_1+b_0D_1]=\tfrac12\ell_2'(x_3).
\]
Both are function elements depending only on $x_3$,
so $\ell_1'$ and $\ell_2'$ are constants.
Since $\ell_1(0)=\ell_2(0)=0$, we obtain
$\ell_1(x_3)=c_1x_3$ and $\ell_2(x_3)=c_2x_3$.
Consequently,
\begin{equation}
\eta=\eta_0(x_1,x_2)+2\varepsilon(x_3)f_3+e(x_3)
+x_3(c_1x_1+c_2x_2),
\qquad c_1,c_2\in\mathbb R.
\label{eq:sectorII-eta}
\end{equation}

\medskip
\noindent\emph{Step 2: Obtain polynomial coefficients in $x_3$.}

Choose $\Lambda=\Lambda(x_3)$ with $\Lambda'=\varepsilon$,
and write $\widetilde P=e^\Lambda Pe^{-\Lambda}$ and
$\widetilde E=e^\Lambda Ee^{-\Lambda}$.
Then $\partial_1\in\widetilde E$ and
$\widetilde T=\partial_3+g(x_1,x_2)\in\widetilde E$.
Set $r=e-\varepsilon^2$. Expanding $\widetilde L_0$ using
\eqref{eq:sectorII-eta} gives
\begin{align}
\widetilde L_0
&=\tfrac12(\partial_1^2+\partial_2^2+\partial_3^2)
-b_0x_1\partial_2-(f_3+\varepsilon)\partial_3\notag\\
&\quad+V_0(x_1,x_2)+x_3V_1(x_1,x_2)
-\tfrac12\bigl(\varepsilon'(x_3)+r(x_3)\bigr).
\label{eq:sectorII-L}
\end{align}
Furthermore,
\[
[\widetilde L_0,\widetilde T]
=(\partial_1g)\partial_1+(\partial_2g)\partial_2
+\varepsilon'(x_3)\partial_3\pmod{\mathcal U_0}.
\]
Transporting this operator back to $E$ and applying
Theorem~\ref{thm:top-coeff} shows that $\varepsilon'$ is
polynomial; if the operator has order zero, the displayed
coefficients vanish. Hence $\varepsilon\in\mathbb R[x_3]$.

Let $d=\deg\varepsilon$, with $d=0$ if $\varepsilon=0$,
and let $m=\max\{2,d\}$.
The commutator formulas give
\[
[g,[g,\widetilde L_0]]
=(\partial_1g)^2+(\partial_2g)^2,
\qquad
[\partial_3,[g,\widetilde L_0]]=0.
\]
Since $[\partial_3,g]=0$, it follows that
\[
\begin{aligned}
\operatorname{ad}_{\widetilde T}^{\,m}(\widetilde L_0)
+\varepsilon^{(m)}\widetilde T
={}&\varepsilon^{(m)}g
+\mathbf 1_{\{m=2\}}
 \bigl((\partial_1g)^2+(\partial_2g)^2\bigr)-\tfrac12r^{(m)}(x_3).
\end{aligned}
\]
Here $\varepsilon^{(m)}$ is constant, so the left-hand side
belongs to $\widetilde E$. It is a multiplication operator,
and therefore is independent of $x_3$. Thus $r^{(m)}$ is
constant and
\[
r\in\mathbb R[x_3],\qquad
\deg r\le\max\{2,d\}\quad(r\ne0).
\]

\medskip
\noindent\emph{Step 3: Exclude $d\ge3$ by increasing weighted-degrees.}

For the remaining steps, consider finite-order differential
operators whose coefficients are smooth in $x_1,x_2$ and
polynomial in $x_3$. Define the weighted degree of a nonzero
term by
\[
\nu\!\left(
a(x_1,x_2)x_3^p
\partial_1^{q_1}\partial_2^{q_2}\partial_3^{q_3}
\right)
=p+q_1+q_2+2q_3,
\]
where $a\in C^\infty(\mathbb R^2)$ and
$p,q_1,q_2,q_3\in\mathbb N_0$.
For a nonzero operator $P$, let $\nu(P)$ be the largest
weighted degree of its nonzero terms in normal form,
and set $\nu(0)=-\infty$.
The product rule gives, for operators $P$ polynomial in
$x_3$ and integers $s\ge0$,
\begin{align*}
\nu([\tfrac12\Delta-f_3\partial_3,P])
&\le\nu(P)+1,\\
\nu([x_3^s\partial_3,P])
&\le\nu(P)+s-1,\\
\nu([x_3^s,P])
&\le\nu(P)+s-3,\\
\nu([-b_0x_1\partial_2+V_0+x_3V_1,P])
&\le\nu(P).
\end{align*}
Consider
\[
S_n=(\operatorname{ad}_{\widetilde L_0})^n(\partial_1)
\in\widetilde E.
\]
Their coefficients are polynomial in $x_3$. Their real span
is finite-dimensional, so a finite basis gives a common
upper bound for $\nu(S_n)$.

Suppose $d\ge3$, and write
$\varepsilon(x_3)=a x_3^d+\cdots$, where $a\ne0$.
Since $\deg r\le d$, only $-a x_3^d\partial_3$ can increase
weighted-degree by $d-1$. Direct calculation gives
\[
S_1=k_1x_1\partial_3+\text{terms of weighted-degree at most }1
\]
and iteratively,
\[
[-a x_3^d\partial_3,x_1x_3^q\partial_3]
=a(d-q)x_1x_3^{q+d-1}\partial_3.
\]
Thus the highest-weighted-degree term of $S_{n+1}$ is
\[
k_1a^n\prod_{j=0}^{n-1}\bigl(d-j(d-1)\bigr)
x_1x_3^{n(d-1)}\partial_3.
\]
Every factor is nonzero because $d/(d-1)$ is not an integer
when $d\ge3$. Its weighted-degree is $n(d-1)+2$, contradicting
the common upper bound. Hence $d\le2$.

\medskip
\noindent\emph{Step 4: Exclude $d\le2$ using the Wronskian lemma.}

Let $a$ be the coefficient of $x_3^2$ in $\varepsilon$.
The preceding bounds give $\nu(S_n)\le n+1$.
For $b(x_1,x_3)\partial_1^i\partial_3^j$ homogeneous of
weighted-degree $n+1$, direct calculation gives
\begin{align*}
[\widetilde L_0,b\partial_1^i\partial_3^j]
&=(\partial_1b)\partial_1^{i+1}\partial_3^j
+(\partial_3b)\partial_1^i\partial_3^{j+1}\\
&\quad
+i k_1x_1b\partial_1^{i-1}\partial_3^{j+1}
+a(2jx_3b-x_3^2\partial_3b)\partial_1^i\partial_3^j\\
&\quad+\text{terms of weighted-degree at most }n+1.
\end{align*}
The term multiplied by $i$ is omitted when $i=0$.
The nonzero displayed terms have weighted-degree $n+2$ and involve
neither $x_2$ nor $\partial_2$.

Starting from $S_0=\partial_1$, the preceding commutator
calculation shows inductively that the terms of weighted
degree $n+1$ in $S_n$ have polynomial coefficients in
$x_1,x_3$ and involve only $\partial_1,\partial_3$.
Write
\[
S_n=\sum_{i,j\ge0}
b_{ij}^{(n)}(x_1,x_3)\partial_1^i\partial_3^j
+\text{terms of weighted degree at most }n,
\]
where the displayed sum consists of the terms of weighted
degree $n+1$. Introduce independent variables $\xi,\zeta$
and define
\[
q_n(x_1,x_3,\xi,\zeta)
=\sum_{i,j\ge0}b_{ij}^{(n)}(x_1,x_3)\xi^i\zeta^j.
\]
Thus $q_n=0$ if and only if the weighted-degree-$(n+1)$
terms of $S_n$ vanish.

For a single term $b\partial_1^i\partial_3^j$,
the terms of weighted-degree $n+2$ in its commutator
with $\widetilde L_0$ are recorded by
\[
\begin{aligned}
&(\partial_1b)\xi^{i+1}\zeta^j
+(\partial_3b)\xi^i\zeta^{j+1}
+i k_1x_1b\xi^{i-1}\zeta^{j+1}
+a\bigl(2j x_3b-x_3^2\partial_3b\bigr)\xi^i\zeta^j.
\end{aligned}
\]
Here the term multiplied by $i$ is omitted when $i=0$.
Summing over $i,j$ gives
\[
\begin{aligned}
q_0&=\xi,\\
q_{n+1}
&=\xi\partial_{x_1}q_n
+(\zeta-a x_3^2)\partial_{x_3}q_n
+k_1x_1\zeta\partial_\xi q_n
+2a x_3\zeta\partial_\zeta q_n.
\end{aligned}
\]
Terms of weighted degree at most $n$ in $S_n$ cannot
contribute to this formula, since commutation with
$\widetilde L_0$ increases weighted degree by at most one.

To interpret this recurrence, consider the auxiliary system
\[
\dot x_1=\xi,\qquad
\dot x_3=\zeta-a x_3^2,\qquad
\dot\xi=k_1x_1\zeta,\qquad
\dot\zeta=2a x_3\zeta.
\]
Its right-hand sides are precisely the four coefficients
in the recurrence. Along any local solution, the chain rule gives
\[
\begin{aligned}
\frac{d}{d\tau}
q_n(x_1(\tau),x_3(\tau),\xi(\tau),\zeta(\tau))
&=\dot x_1\partial_{x_1}q_n
+\dot x_3\partial_{x_3}q_n
+\dot\xi\partial_\xi q_n
+\dot\zeta\partial_\zeta q_n\\
&=q_{n+1}(x_1(\tau),x_3(\tau),\xi(\tau),\zeta(\tau)),
\end{aligned}
\]
where the partial derivatives are evaluated along the solution.
Since $q_0=\xi$, induction yields
\[
q_n(x_1(\tau),x_3(\tau),\xi(\tau),\zeta(\tau))
=\frac{d^n\xi}{d\tau^n}(\tau).
\]

Suppose now that $q_N=0$ as a polynomial for some $N\ge1$.
Then $\xi^{(N)}=0$ along every local solution.
Consequently, $\xi(\tau)$ is polynomial in $\tau$,
and so is $x_1(\tau)$ because $\dot x_1=\xi$.

The equations for $x_3,\zeta$ are independent of $x_1,\xi$.
Fix $x_3(0)=0$ and $\zeta(0)=1$.
By uniqueness, the two solutions with
\[
(x_1(0),\xi(0))=(1,0)
\qquad\text{and}\qquad
(x_1(0),\xi(0))=(0,1)
\]
have the same functions $x_3(\tau),\zeta(\tau)$.
Their $x_1$ components therefore give polynomial solutions
$p,q$ of the same equation
\[
u''=k_1\zeta(\tau)u,
\]
with
\(
p(0)=1,p'(0)=0,
q(0)=0,q'(0)=1.
\)
Their Wronskian has derivative zero and initial value one,
so $pq'-p'q=1$ on a common interval, hence as a polynomial
identity. Lemma~\ref{lem:constant-wronskian} implies that
their span contains the constant function $1$.
Substitution into the equation gives $k_1\zeta(\tau)=0$,
contrary to $k_1\ne0$ and $\zeta(0)=1$.

Therefore $q_n\ne0$ for every $n$, and
$\nu(S_n)=n+1$, again contradicting the common upper bound.
The assumed nonzero coefficient matrix is therefore impossible,
and $k_1=k_2=h_2=0$ in the original coordinates.
\end{proof}






\begin{proof}[Proof of Theorem~\ref{thm:main}]
By Theorem~\ref{thm:shi-yau-affine}, the Wong matrix is affine, so it
suffices to show that every one of its coefficients vanishes. This is
done in three steps.
\begin{enumerate}
    \item Theorem~\ref{thm:w12-constant} gives $b_1=b_2=0$, i.e., $\omega_{12}=b_0$ is constant.
    \item Theorem~\ref{thm:sector-I-smooth} gives $k_3=h_3=0$, i.e., $\omega_{13}$ and $\omega_{23}$ are independent of $x_3$.
    \item Theorem~\ref{thm:remaining-visible} gives $k_1=k_2=h_2=0$, and Theorem~\ref{thm:bianchi-affine} gives $h_1=k_2=0$; hence $\omega_{13}$ and $\omega_{23}$ are constant.
\end{enumerate}
All coefficients of the affine Wong matrix have been eliminated, so
$\Omega$ is constant.
\end{proof}
\appendix
\section{Proof of the Kernel of Euler Operator}\label{app:euler}
\begin{lemma}[Global smooth Euler equation]\label{lem:global-smooth-euler}
Let $Q\in\mathbb C[s]\setminus\{0\}$ and let $A=t\partial_t$. Then
\[
 \ker_{C^\infty(\mathbb R;\mathbb C)}Q(A)
 =\operatorname{span}_{\mathbb C}
   \{t^n:n\in\mathbb N_0,\ Q(n)=0\}.
\]
In particular, every globally smooth solution of a nonzero polynomial Euler equation is a
finite polynomial.
\end{lemma}

\begin{proof}
If $Q$ is a nonzero constant, both sides of the asserted identity are zero.
Assume henceforth that $N:=\deg Q\ge1$.
For $u\in C^\infty(\mathbb R;\mathbb C)$ and every $n\in\mathbb N_0$, the Leibniz rule gives
\[
 \bigl(Q(A)u\bigr)^{(n)}(0)=Q(n)u^{(n)}(0).
\]
Suppose that $Q(A)u=0$, and set
\[
 S=\{n\in\mathbb N_0:Q(n)=0\},
 \qquad
 p(t)=\sum_{n\in S}\frac{u^{(n)}(0)}{n!}t^n.
\]
The set $S$ is finite, $Q(A)p=0$, and $v=u-p$ satisfies $Q(A)v=0$ and is
flat at $t=0$.

We now rigorously exclude nonzero flat solutions. On $t>0$, set $t=e^s$
and $V_+(s)=v(e^s)$; on $t<0$, set $t=-e^s$ and
$V_-(s)=v(-e^s)$. On both half-lines, $A=\partial_s$, and hence
\[
 Q(\partial_s)V_\pm=0.
\]
For either choice of sign, put
\[
W_\pm(s)=\bigl(V_\pm(s),V_\pm'(s),\ldots,V_\pm^{(N-1)}(s)\bigr)^T.
\]
After dividing the equation by the nonzero leading coefficient of $Q$,
there is a constant complex companion matrix $M$ such that
$W_\pm'=MW_\pm$. Flatness of $v$ and the chain rule imply
\[
\|W_\pm(s)\|=o(e^{Ls})\qquad(s\to-\infty)
\quad\hbox{for every }L>0.
\]
Choose a Euclidean norm on $\mathbb C^N$ and its induced matrix norm.
Fix $s_0\in\mathbb R$. Uniqueness for this constant-coefficient system gives,
for $s\le s_0$,
\[
W_\pm(s_0)=e^{(s_0-s)M}W_\pm(s),
\qquad
\|W_\pm(s_0)\|
\le e^{(s_0-s)\|M\|}\|W_\pm(s)\|.
\]
Choose $L>\|M\|$ and let $s\to-\infty$. The right-hand side tends to zero,
so $W_\pm(s_0)=0$. Since $s_0$ was arbitrary, $V_+=V_-=0$.
Consequently $v=0$ on both half-lines and at the origin, proving the claim.
\end{proof}


\end{revision}

\section*{Acknowledgment}
The author thanks his close collaborator Professor Zeju Sun for many
helpful suggestions, and Tirant Cai for her constant support.

\bibliographystyle{plainnat}
\bibliography{ref}

\end{document}
